% -*- coding: utf-8 -*-
\chapter{研究発表スライドの作成}
\label{chap5}

\begin{chapterbox}
\textbf{【この章で学ぶこと】}
\begin{itemize}
\item \textbf{AIを活用して論理的で説得力のある研究発表スライドを作成できる} --- 発表の構成案からスライド内容まで、AIを効果的に活用した作成プロセスを実践できる
\item \textbf{デザインAIツールを使った効率的なスライド作成法を習得する} --- GammaやCanva AIなどのツールを活用し、視覚的に優れたスライドを短時間で作成できる
\item \textbf{図表・概念図の生成と批判的な検討ができる} --- AIによる図表生成の可能性と限界を理解し、適切に活用・検証できる
\end{itemize}
\end{chapterbox}
\vspace{5mm}

%======================================================================
\section{研究発表スライドの基本}
\label{sec:5.1}
\index{けんきゅうはっぴょうすらいど@研究発表スライド}

\subsection{学会発表の基本構成}
\index{がっかいはっぴょう@学会発表}

研究発表のスライドは、論文と同様に論理的な構成が求められます。ただし、スライドは「読む」ものではなく「見せる」ものであるため、論文とは異なるアプローチが必要です。

\textbf{標準的な発表構成:}

\begin{table}[htbp]
\centering
\caption{研究発表スライドの標準的な構成}
\label{tab:slide-structure}
\begin{tabular}{lll}
\toprule
スライド & 内容 & 目的 \\
\midrule
タイトル & 研究タイトル、著者名、所属 & 研究の第一印象を決める \\
背景 & 研究分野の現状と課題 & 聴衆に問題の重要性を理解させる \\
目的 & 本研究で解決する問題 & 研究のゴールを明確にする \\
関連研究 & 先行研究と本研究の位置づけ & 新規性の根拠を示す \\
提案手法 & アプローチの概要と詳細 & 技術的な貢献を説明する \\
実験設定 & データセット、評価指標、比較手法 & 実験の妥当性を示す \\
結果 & 主要な実験結果 & 提案手法の有効性を示す \\
考察 & 結果の解釈、分析 & 結果の意味を深く理解させる \\
まとめ & 貢献の要約と将来展望 & 研究のインパクトを印象づける \\
\bottomrule
\end{tabular}
\end{table}

\subsection{時間配分の目安（10分発表の場合）}

10分間の口頭発表は、学会で最も一般的な発表形式の一つです。以下に推奨される時間配分を示します。

\begin{table}[htbp]
\centering
\caption{10分発表の推奨時間配分}
\label{tab:time-allocation}
\begin{tabular}{lll}
\toprule
セクション & 時間 & スライド枚数（目安） \\
\midrule
タイトル・自己紹介 & 0.5分 & 1枚 \\
背景・問題設定 & 2分 & 2--3枚 \\
提案手法 & 3分 & 3--4枚 \\
実験結果 & 2.5分 & 2--3枚 \\
考察・まとめ & 2分 & 2枚 \\
\midrule
\textbf{合計} & \textbf{10分} & \textbf{10--13枚} \\
\bottomrule
\end{tabular}
\end{table}

1枚のスライドに費やす時間は、平均して45秒から1分程度です。情報を詰め込みすぎると、聴衆は内容を理解する前に次のスライドに移ってしまいます。

\subsection{良いスライドの原則}

\begin{itembox}[l]{1スライド1メッセージの原則}
各スライドが伝えるべきメッセージは1つだけです。複数のメッセージを1枚に詰め込むと、聴衆はどこに注目すべきかわからなくなります。スライドのタイトルに、そのスライドのメッセージを端的に表現しましょう。

\textbf{悪い例のタイトル:} 「実験結果」

\textbf{良い例のタイトル:} 「提案手法は全データセットでベースラインを上回った」
\end{itembox}

\textbf{文字量の制限:}

\begin{itemize}
\item タイトル: 1行（最大15語程度）
\item 本文: 1スライドあたり最大6行
\item フォントサイズ: タイトル28pt以上、本文20pt以上
\item 箇条書きは3--5項目まで
\end{itemize}

\textbf{視覚的な原則:}

\begin{itemize}
\item 図表を積極的に使う（テキストより図の方が記憶に残る）
\item 色の使いすぎを避ける（メインカラー + アクセントカラーの2--3色）
\item 余白を十分にとる（スライドの30\%以上は空白にする）
\item アニメーションは必要最小限にする
\end{itemize}

\begin{itembox}[l]{Tip: スライド作成の出発点}
スライドを作成する前に、まず「各スライドで伝えたいメッセージ」を箇条書きで書き出しましょう。メッセージの一覧が論理的に流れていれば、スライド作成はスムーズに進みます。この段階でAIを活用すると効率的です。
\end{itembox}

%======================================================================
\section{AIによる構成案の生成}
\label{sec:5.2}
\index{こうせいあん@構成案}

\subsection{研究テーマの構造化}
\index{えれべーたーぴっち@エレベーターピッチ}

AIに構成案を依頼する前に、自分の研究を1--2文で簡潔に表現できるようにしましょう。これは「エレベーターピッチ」とも呼ばれ、短時間で研究の本質を伝える技術です。

\begin{promptbox}
私の研究を以下の情報をもとに、1-2文の簡潔な説明にまとめてください。

- 研究分野: [分野]
- 解決する問題: [問題]
- 提案するアプローチ: [手法の概要]
- 主な結果: [成果]
- 研究の意義: [意義]

対象聴衆: [例: 同じ分野の研究者 / 異分野の研究者 / 一般的な理系の聴衆]
\end{promptbox}

\subsection{AIに構成案を提案させるプロンプト設計}

研究テーマが整理できたら、AIに発表スライドの構成案を提案させます。

\begin{promptbox}
以下の研究について、10分間の学会口頭発表のスライド構成案を
作成してください。

【研究概要】
タイトル: [研究タイトル]
分野: [研究分野]
概要: [1-2文の研究概要]

【詳細情報】
背景: [研究の背景と動機]
提案手法: [手法の概要]
主要結果: [主な実験結果]
結論: [研究の結論]

【発表条件】
発表時間: 10分（質疑5分）
対象聴衆: [聴衆の専門性レベル]
発表言語: [日本語/英語]

以下の形式で構成案を作成してください:
- スライド番号
- スライドタイトル（メッセージ型）
- 含めるべき内容（箇条書き3-5項目）
- 想定所要時間
- 視覚要素の提案（図表、グラフなど）
\end{promptbox}

\subsection{生成された構成案の批判的検討}

AIが生成した構成案は、必ず以下の4つの観点から批判的に検討します。

\textbf{1. 論理一貫性の確認:}

スライドのタイトルだけを順番に読んで、研究のストーリーが自然に流れているか確認します。

\begin{promptbox}
以下のスライド構成案について、論理的な一貫性を確認してください。
スライドタイトルの順序を読むだけで、研究のストーリーが
伝わるかどうかを評価してください。

論理の飛躍、不足している要素、順序の改善案があれば指摘してください。

[スライドタイトルのリスト]
\end{promptbox}

\textbf{2. 網羅性の確認:}

必要な情報がすべて含まれているか確認します。特に以下の要素は必須です。

\begin{itemize}
\item 研究の動機（なぜこの研究が重要か）
\item 先行研究との差別化（何が新しいのか）
\item 手法の核心的なアイデア（どのように解決するか）
\item 定量的な結果（どの程度うまくいったか）
\item 限界と将来展望（残された課題は何か）
\end{itemize}

\textbf{3. 独自性の確認:}

構成案が一般的すぎないか、自分の研究の独自性が反映されているか確認します。AIは汎用的な構成を提案する傾向があるため、自分の研究の強みが際立つように修正を加えます。

\textbf{4. 時間配分の確認:}

各スライドの想定所要時間を合計し、発表時間内に収まるか確認します。実際に声に出してリハーサルすると、想定よりも時間がかかることが多いため、余裕を持った時間配分にしましょう。

\begin{itembox}[l]{Tip: 早い段階でのフィードバック}
構成案の段階で指導教員やゼミのメンバーにフィードバックをもらうのが効果的です。スライドを作り込んでからの大幅な修正は時間のロスが大きいため、早い段階で方向性を固めましょう。
\end{itembox}

%======================================================================
\section{各スライドの内容生成}
\label{sec:5.3}

\subsection{研究背景の文章生成}

研究背景のスライドは、聴衆の注意を引き、研究の重要性を理解してもらうための出発点です。

\begin{promptbox}
以下の研究テーマの背景を、学会発表スライド用の簡潔な文章として
生成してください。

【研究テーマ】: [テーマ]
【解決したい問題】: [問題]
【対象聴衆】: [聴衆の属性]

以下の条件を守ってください:
1. 箇条書き形式で、1項目あたり1行（15-20文字程度）
2. 専門用語は[聴衆の属性]が理解できるレベルにする
3. 問題の重要性が直感的に伝わるデータや事実を含める
4. 最後に「この問題をどう解決するか？」への橋渡しとなる文で締める
5. 合計5-6行以内
\end{promptbox}

\subsection{研究手法の説明文生成}

提案手法の説明は、発表の核心部分です。複雑な手法を聴衆に理解してもらうためには、段階的な説明が重要です。

\begin{promptbox}
以下の研究手法を、学会発表スライド用に段階的に説明する文章を
生成してください。

【手法の概要】: [手法の詳細な説明]

以下の構成で、3枚のスライドに分けてください:

スライド1: 手法の全体像（Overview）
- 手法の核心的なアイデアを1文で
- 全体のパイプラインを箇条書き3-4項目で

スライド2: 技術的な詳細（Key Technique）
- 最も重要な技術的貢献に焦点
- 数式は最小限（1-2個まで）
- 直感的な説明を優先

スライド3: 実装のポイント
- 実用的な工夫
- 計算コストや効率性

各スライドは5行以内で、図表の提案も含めてください。
\end{promptbox}

\subsection{研究結果のまとめ生成}

\begin{promptbox}
以下の実験結果を、学会発表スライド用にまとめてください。

【実験結果】:
[結果のデータを貼り付け]

以下の形式でお願いします:
1. 主要な発見を重要度順に3つ挙げる
2. 各発見を1行で端的に表現する
3. 数値は効果的に強調する（例: "精度12.5\%向上"）
4. 図表（グラフの種類）の提案
5. 聴衆が「なるほど」と思えるような解釈を1文添える
\end{promptbox}

\subsection{専門用語と平易な表現のバランス調整}

学会発表では、聴衆の専門性に応じて表現のレベルを調整する必要があります。

\begin{promptbox}
以下のスライド内容について、対象聴衆に合わせて表現を調整してください。

【現在の内容】:
[スライドの内容]

【対象聴衆】: [例: 情報科学を専攻する修士学生]

調整の方針:
1. 聴衆が確実に知っている用語はそのまま使用
2. 聴衆が知らない可能性のある用語は、初出時に簡単な説明を付加
3. 略語は初出でフルスペルを示す
4. 過度に専門的な詳細は省略し、直感的な説明に置き換える
\end{promptbox}

\subsection{AIの出力を自分の言葉に修正する重要性}

AIが生成したスライド内容は、あくまで出発点です。以下の理由から、必ず自分の言葉に修正してください。

\begin{itembox}[l]{なぜ自分の言葉に修正が必要か}
\begin{enumerate}
\item \textbf{発表の自然さ:} AIが生成した文章をそのまま読み上げると、不自然な発表になります。自分が普段使う言葉遣いに修正することで、自信を持って発表できます。
\item \textbf{質疑への対応:} 自分の言葉でスライドを作成していれば、質疑応答で深い質問をされても自然に回答できます。AIの文章をそのまま使っていると、内容を十分に理解していないことが露呈するリスクがあります。
\item \textbf{独自の視点の反映:} AIは汎用的な表現を生成しますが、あなた自身の研究に対する独自の洞察や情熱は、自分の言葉でしか表現できません。
\end{enumerate}
\end{itembox}

\begin{itembox}[l]{Tip: 違和感を大切にする}
AIが生成した内容を読んで「これは自分の考えと違う」と感じた箇所は、積極的に修正しましょう。その違和感こそが、あなた自身の研究に対する理解の深さを反映しています。
\end{itembox}

%======================================================================
\section{デザインAIによるスライド作成}
\label{sec:5.4}
\index{でざいんAI@デザインAI}

\subsection{Gamma --- テキストからスライドを自動生成}
\index{Gamma@Gamma}

Gamma（gamma.app）は、テキスト入力から高品質なプレゼンテーションスライドを自動生成するAIツールです。研究発表スライドの初稿を短時間で作成するのに適しています。

\textbf{基本的な使い方:}

\begin{enumerate}
\item Gammaのウェブサイトにアクセスし、アカウントを作成する
\item 「Create new」から「Presentation」を選択する
\item 研究の概要テキストを入力するか、アウトライン形式で構成を指定する
\item AIがスライドのデザインと内容を自動生成する
\item 生成されたスライドを編集・カスタマイズする
\end{enumerate}

\textbf{研究発表に使う際のプロンプト例:}

\begin{promptbox}
Create a 12-slide academic research presentation:

Title: [研究タイトル]

Slide 1: Title slide
Slide 2: Research Background - [背景の要点]
Slide 3: Problem Statement - [問題の記述]
Slide 4: Related Work - [先行研究の位置づけ]
Slide 5-7: Proposed Method - [手法の説明]
Slide 8: Experimental Setup - [実験設定]
Slide 9-10: Results - [結果の要点]
Slide 11: Discussion - [考察の要点]
Slide 12: Conclusion and Future Work - [結論]

Style: Clean, minimal, academic
Color scheme: Blue and white
\end{promptbox}

\textbf{Gammaの長所:}
\begin{itemize}
\item テキストから短時間でスライドを生成できる
\item デザインが洗練されている
\item 図表のプレースホルダーが自動配置される
\item エクスポート形式が豊富（PDF, PPTX等）
\end{itemize}

\textbf{Gammaの短所:}
\begin{itemize}
\item 細かいレイアウトの調整には限界がある
\item 数式の表示には対応していない場合がある
\item 学術発表特有のフォーマット（ヘッダー/フッターなど）のカスタマイズが限定的
\item 無料プランでは一部機能が制限される
\end{itemize}

\subsection{Canva AI --- テンプレートベースのデザイン支援}
\index{Canva AI@Canva AI}

Canva（canva.com）は、テンプレートベースのデザインツールで、AI機能（Magic Design）によるスライド作成支援を提供しています。

\textbf{基本的な使い方:}

\begin{enumerate}
\item Canvaにアクセスし、「プレゼンテーション」を選択する
\item テンプレートを選択する（学術発表向けテンプレートが複数用意されている）
\item Magic Designにテーマを入力すると、テンプレートに合わせたスライドが提案される
\item 各スライドの内容を編集する
\item 図表やアイコンをCanvaの素材ライブラリから追加する
\end{enumerate}

\textbf{学術発表に適したテンプレートの選び方:}

\begin{itemize}
\item シンプルで余白が十分にあるもの
\item 色数が少ないもの（2--3色）
\item フォントが読みやすいもの（サンセリフ体推奨）
\item 装飾的な要素が少ないもの
\item グラフや表のプレースホルダーがあるもの
\end{itemize}

\begin{itembox}[l]{Tip: 最終調整はPowerPointで}
デザインAIで作成したスライドは、必ずPowerPointやKeynote形式でエクスポートし、最終的な微調整はこれらのツールで行いましょう。特に数式、特殊な図表、アニメーションについては、PowerPoint等で仕上げた方が柔軟性が高いです。
\end{itembox}

\subsection{デザインAIの使い分け}

\begin{table}[htbp]
\centering
\caption{デザインAIツールの比較}
\label{tab:design-ai-comparison}
\begin{tabular}{lll}
\toprule
項目 & Gamma & Canva AI \\
\midrule
強み & テキストからの自動生成 & テンプレートの豊富さ \\
作成速度 & 非常に速い & 速い \\
カスタマイズ性 & 中程度 & 高い \\
数式対応 & 限定的 & 限定的 \\
共同編集 & 対応 & 対応 \\
学術向けテンプレート & 少なめ & 多い \\
推奨用途 & 初稿の高速作成 & デザインの洗練 \\
\bottomrule
\end{tabular}
\end{table}

%======================================================================
\section{図表・概念図の生成}
\label{sec:5.5}
\index{がいねんず@概念図}

\subsection{画像生成AIによる概念図}
\index{がぞうせいせいAI@画像生成AI}

研究発表では、手法の概念図やシステムの全体像を示す図が非常に重要です。画像生成AI（DALL-E、Stable Diffusion等）を活用して、概念図のアイデアを得ることができます。

\textbf{活用できる場面:}

\begin{itemize}
\item 研究のイメージを伝えるための概要図
\item 応用分野のイラスト（医療、環境、ロボティクスなど）
\item スライドの背景画像やアイコン
\end{itemize}

\textbf{活用が難しい場面:}

\begin{itemize}
\item 正確な技術的アーキテクチャ図（ネットワーク構造など）
\item 数値データを含むグラフや表
\item 数式や記号を含む図
\end{itemize}

\textbf{プロンプト例:}

\begin{promptbox}
Create a clean, minimalist scientific illustration showing the concept
of transfer learning in natural language processing.
Show knowledge flowing from a large pre-trained model to a smaller
task-specific model. Use a blue and white color scheme.
Style: technical diagram, flat design, no text.
\end{promptbox}

\begin{itembox}[l]{Tip: 2段階アプローチ}
画像生成AIは概念的なイラストの「アイデア出し」には有用ですが、最終的な研究発表用の図は、専用のツール（PowerPoint、Illustrator、draw.io等）で自作することを推奨します。生成画像には不正確な要素が含まれている可能性があり、学術発表での使用には注意が必要です。
\end{itembox}

\subsection{テキストベース図表ツール --- Mermaid.jsの活用}
\index{Mermaid.js@Mermaid.js}

Mermaid.jsは、テキストベースで図表を生成するツールです。AIにMermaidコードを生成させることで、フローチャート、シーケンス図、状態遷移図などを効率的に作成できます。

\textbf{フローチャートの生成プロンプト:}

\begin{promptbox}
以下の研究手法のパイプラインを、Mermaid.jsのフローチャートとして
表現してください。

【手法のパイプライン】:
1. 入力データの前処理（テキストのトークナイズ、正規化）
2. 特徴抽出（BERTによるエンコーディング）
3. 分類層（全結合層 + Softmax）
4. 出力（カテゴリラベル）

各ステップ間のデータの流れが明確にわかるようにしてください。
重要なステップは色を変えて強調してください。
\end{promptbox}

\textbf{AIが生成するMermaidコードの例:}

\begin{lstlisting}[language={}]
graph TD
    A[入力テキスト] --> B[前処理]
    B --> B1[トークナイズ]
    B --> B2[正規化]
    B1 --> C[特徴抽出]
    B2 --> C
    C --> C1[BERT エンコーダ]
    C1 --> D[分類層]
    D --> D1[全結合層]
    D1 --> D2[Softmax]
    D2 --> E[カテゴリラベル出力]

    style C1 fill:#4A90D9,stroke:#333,color:#fff
    style D2 fill:#4A90D9,stroke:#333,color:#fff
\end{lstlisting}

\textbf{シーケンス図の生成プロンプト:}

\begin{promptbox}
以下のシステムの処理フローを、Mermaid.jsのシーケンス図として
表現してください。

【処理フロー】:
1. ユーザーがクエリを入力
2. フロントエンドがAPIサーバーにリクエストを送信
3. APIサーバーがデータベースに問い合わせ
4. データベースが結果を返す
5. APIサーバーがAIモデルに推論を依頼
6. AIモデルが結果を返す
7. APIサーバーがフロントエンドにレスポンスを返す
8. フロントエンドがユーザーに結果を表示
\end{promptbox}

\textbf{研究手法の比較図:}

\begin{promptbox}
以下の3つの手法の比較を、Mermaid.jsの図として表現してください。
各手法の入力、処理ステップ、出力を並列に表示し、
共通点と相違点が一目でわかるようにしてください。

手法A: [概要]
手法B: [概要]
手法C（提案手法）: [概要]
\end{promptbox}

\textbf{Mermaidコードの活用方法:}

生成されたMermaidコードは、以下の方法でスライドに組み込むことができます。

\begin{enumerate}
\item \textbf{Mermaid Live Editor}（mermaid.live）でSVGまたはPNG画像としてエクスポート
\item \textbf{VS Code拡張機能}（Markdown Preview Mermaid Support）でプレビュー
\item \textbf{GitHubのMarkdown}はMermaidを直接レンダリング可能
\item エクスポートした画像をPowerPointやKeynoteに貼り付け
\end{enumerate}

\subsection{生成画像の著作権と正確性の検証}
\index{ちょさくけん@著作権}

AIで生成した図表をスライドに使用する際は、以下の点に注意が必要です。

\textbf{著作権に関する注意:}

\begin{itemize}
\item 画像生成AIの利用規約を確認する（商用利用、学術利用の可否）
\item DALL-Eで生成した画像はOpenAIの利用規約に基づいて使用可能
\item Stable Diffusionはモデルのライセンスによって異なる
\item 学会の規定によっては、AI生成画像の使用に制限がある場合がある
\end{itemize}

\textbf{正確性の検証:}

\begin{itemize}
\item AIが生成した概念図の要素が正確かどうかを確認する
\item 生成画像に含まれるテキストは不正確なことが多いため、自分で修正する
\item 技術的な図（ネットワーク構造、数式を含む図）は特に注意して検証する
\item 不正確な図を使用すると、査読者や聴衆の信頼を損なう可能性がある
\end{itemize}

\begin{itembox}[l]{Tip: Overview図の作成}
学会発表では、手法の全体像を示す「Overview図」が特に重要です。AIで概念図のラフ案を生成し、それを参考にしてPowerPointやdraw.ioで正確な図を自作するという2段階のアプローチが効果的です。最終的な図には、自分で正確性を保証できるものだけを使用しましょう。
\end{itembox}

%======================================================================
\section{相互レビューと改善}
\label{sec:5.6}
\index{そうごれびゅー@相互レビュー}

\subsection{ピアレビューのチェックポイント}

スライドが完成したら、発表前にピアレビュー（相互レビュー）を行いましょう。以下のチェックポイントを活用してください。

\textbf{内容面のチェック:}

\begin{itemize}
\item 研究の動機が明確に伝わるか
\item 提案手法の核心的なアイデアが理解できるか
\item 実験結果が説得力を持って提示されているか
\item 結論が実験結果に基づいているか
\item スライド間の論理的なつながりがスムーズか
\end{itemize}

\textbf{デザイン面のチェック:}

\begin{itemize}
\item テキストが読みやすいか（フォントサイズ、コントラスト）
\item 図表が鮮明で、説明なしでも理解できるか
\item 色の使い方が一貫しているか
\item 余白が十分にあるか
\item アニメーションが適切か（過度でないか）
\end{itemize}

\textbf{発表面のチェック:}

\begin{itemize}
\item 制限時間内に収まるか
\item 各スライドのメッセージが明確か
\item 専門用語が適切に説明されているか
\item 質疑応答で予想される質問への準備ができているか
\end{itemize}

\subsection{AIに「聴衆として」フィードバックさせる}

人間のレビューに加えて、AIにも聴衆の視点からフィードバックを求めることができます。

\textbf{基本プロンプト:}

\begin{promptbox}
あなたは[分野名]の学会に参加している研究者です。
以下のスライド内容（テキスト版）を読み、発表に対する
フィードバックを提供してください。

【スライド内容】
スライド1: [タイトルと内容]
スライド2: [タイトルと内容]
...

以下の観点からフィードバックをお願いします:
1. ストーリーの流れは自然か？論理的な飛躍はないか？
2. 背景と動機は十分に説得力があるか？
3. 提案手法の説明で不明確な点はないか？
4. 実験結果の提示は効果的か？
5. 聴衆として、質疑応答で聞きたい質問は何か？（3つ挙げてください）
6. 全体的な改善提案
\end{promptbox}

\textbf{異なる聴衆レベルでのフィードバック:}

\begin{promptbox}
以下のスライド内容について、異なる立場の聴衆からのフィードバックを
シミュレーションしてください。

聴衆A: この分野の専門家（教授レベル）
聴衆B: 同じ学科の別分野の大学院生
聴衆C: 理系だが異なる学部の学部生

それぞれの聴衆が「理解できなかった」と感じる可能性が
高いポイントを指摘してください。

[スライド内容]
\end{promptbox}

\textbf{質疑応答の準備:}

\begin{promptbox}
以下の研究発表スライドの内容をもとに、学会の質疑応答で
想定される質問を10個生成してください。

質問は以下のカテゴリに分類してください:
1. 手法に関する技術的な質問（3問）
2. 実験設定に関する質問（2問）
3. 結果の解釈に関する質問（2問）
4. 研究の限界や将来展望に関する質問（2問）
5. 厳しい指摘・批判的な質問（1問）

各質問に対する推奨回答の骨子も示してください。

[スライド内容]
\end{promptbox}

\begin{itembox}[l]{Tip: AIと人間のフィードバックの組み合わせ}
AIからのフィードバックと人間からのフィードバックは、異なる強みを持っています。AIは論理的な一貫性や網羅性のチェックに優れていますが、発表の「温度感」や「説得力」は人間のフィードバックの方が的確です。両方を組み合わせることで、より完成度の高い発表になります。
\end{itembox}

%======================================================================
\section*{演習問題}

\begin{enumerate}

\item \textbf{演習1: 1スライド1メッセージの実践}

自分の研究（または興味のある研究トピック）について、以下を行いなさい。
\begin{enumerate}
\item 10分間の口頭発表を想定し、各スライドで伝えたいメッセージを箇条書きで12個書き出す
\item 各メッセージが「1スライド1メッセージ」の原則を満たしているか確認する（複数のメッセージが混在していたら分割する）
\item メッセージの順序を読むだけで、研究のストーリーが自然に流れるか確認する
\item 必要に応じてメッセージの追加・削除・順序変更を行い、最終的な構成案を完成させる
\end{enumerate}

\item \textbf{演習2: AIによる構成案の比較}

本章のプロンプトテンプレートを使って、2つの異なるAI（例: Claude と ChatGPT）にスライド構成案を生成させ、以下を行いなさい。
\begin{enumerate}
\item 2つの構成案の共通点と相違点を整理する
\item 各構成案の長所と短所を分析する
\item 2つの構成案の良い部分を組み合わせて、最終的な構成案を作成する
\item 最終構成案が、本章で述べた4つの検討観点（論理一貫性、網羅性、独自性、時間配分）を満たしているか確認する
\end{enumerate}

\item \textbf{演習3: Mermaid.jsによる図表作成}

自分の研究手法のパイプライン（または興味のある手法）について、以下を行いなさい。
\begin{enumerate}
\item AIにMermaid.jsのフローチャートコードを生成させる
\item Mermaid Live Editor（mermaid.live）でコードを貼り付け、図を確認する
\item 図が正確に手法を表現しているか検証し、必要に応じてコードを修正する
\item SVGまたはPNG形式でエクスポートし、スライドに挿入する
\item テキストベース図表ツールの長所と短所をまとめる
\end{enumerate}

\item \textbf{演習4: デザインAIの活用}

演習1または2で作成した構成案をもとに、以下を行いなさい。
\begin{enumerate}
\item Gamma（gamma.app）を使ってスライドの初稿を自動生成する
\item 生成されたスライドのデザインと内容を評価する（良い点、改善すべき点）
\item 内容を自分の言葉に修正する
\item デザインの改善（色、フォント、レイアウト、余白）を行う
\item PowerPointまたはKeynote形式でエクスポートし、最終的な微調整を行う
\end{enumerate}

\item \textbf{演習5: AIを活用した質疑応答準備}

完成したスライド（演習4の成果物、またはこれまでに作成したスライド）について、以下を行いなさい。
\begin{enumerate}
\item 本章のプロンプトを使って、AIに想定質問を10個生成させる
\item 各質問に対する回答を自分で準備する（AIの推奨回答を参考にしつつ、自分の言葉で）
\item 特に難しいと感じた質問3つについて、回答を声に出して練習する
\item ゼミや研究室メンバーの前でリハーサル発表を行い、実際にどのような質問が出るかを記録する
\item AIが予測した質問と実際の質問を比較し、AIによる質疑準備の有効性を評価する
\end{enumerate}

\end{enumerate}

\vspace{10mm}
\begin{itembox}[l]{本章のまとめ}
研究発表スライドの作成は、研究内容の理解と伝達力の両方が求められる総合的な作業です。AIは構成案の生成、内容の整理、図表の作成、フィードバックの提供など、多くの段階で強力な支援を提供してくれます。しかし、最も重要なのは「自分の研究を自分の言葉で伝える」ことです。AIの出力は必ず自分で検証・修正し、発表のリハーサルを十分に行いましょう。良い発表は、優れたスライドと十分な練習の両方から生まれます。
\end{itembox}
